\documentclass[conference]{IEEEtran}
\usepackage{amsmath}
\usepackage{balance}
\begin{document}
\title{Exploring Integrated and Decision-Weighted Gradients for Image Attribution}
\author{\IEEEauthorblockN{Taher Akolawala}}
\maketitle
\begin{abstract}
We implemented two path-based attribution methods in PyTorch for a pretrained ImageNet ResNet-50: Integrated Gradients (IG) and an Integrated Decision Gradients (IDG)-inspired, decision-weighted variant. Both trace a straight path from a zero tensor to a normalized input image and use input gradients to assign pixel importance. The decision-weighted method allocates samples according to changes in the target logit and weights gradients by the directional derivative along the path. In a macaque example, its heatmap appeared more concentrated on the animal than the IG heatmap. This observation is qualitative; the examples were not produced under a controlled same-image benchmark, and we did not measure faithfulness or stability. The experiment therefore illustrates a visual difference without establishing that either method provides a better explanation.
\end{abstract}
\begin{IEEEkeywords}
explainable AI, integrated gradients, image attribution, ResNet-50, gradient saturation
\end{IEEEkeywords}

\section{Introduction}
Attribution methods assign scores to input pixels to help examine why an image classifier selected a class. A heatmap can support model inspection, but visual sharpness alone is not proof that an explanation is faithful. Integrated Gradients (IG) integrates the model's input gradient along a path from a reference input to the image and was introduced with axiomatic motivation \cite{ig}. Integrated Decision Gradients (IDG) addresses path regions where the output changes little, known as saturation, by concentrating attention near consequential score changes \cite{idg}. We implemented IG and a decision-weighted variant directly in PyTorch, without an attribution library.

Our aim was to examine how path sampling and decision-sensitive weighting affect image attributions while keeping the model, preprocessing, and target-score definition explicit. Baseline selection and numerical integration can alter the resulting maps, so we distinguish qualitative observations from conclusions that would require controlled evaluation.

\section{Common Experimental Setup}
For both methods, we loaded the torchvision ResNet-50 weights named \texttt{IMAGENET1K\_V1}, set the model to evaluation mode, resized the image to 256 pixels, took a 224$\times$224 center crop, converted it to a tensor, and applied ImageNet channel normalization. We used the predicted top class as the attribution target. The baseline was a zero tensor in the \emph{normalized} input space. Thus ``zero baseline'' here means zero after normalization; it should not automatically be interpreted as a black RGB image.

Let $F_c(x)$ be the selected class logit, $x'$ the baseline, $x$ the normalized input, and $\Delta=x-x'$. The shared path is $\gamma(\alpha)=x'+\alpha\Delta$ for $0\leq\alpha\leq1$. We computed attribution for the target logit rather than a softmax probability. For visualization, we summed the absolute attribution over the three color channels at each pixel. Absolute-value aggregation removes sign, so the heatmaps do not distinguish evidence supporting the class from evidence opposing it.

\section{Integrated Gradients Implementation}
For input coordinate $i$, continuous IG is
\begin{equation}
\mathrm{IG}_i(x)=\Delta_i\int_0^1\frac{\partial F_c(\gamma(\alpha))}{\partial x_i}\,d\alpha.
\end{equation}
We sampled 50 uniformly spaced path positions, including both endpoints, obtained the input gradient at each position by backpropagation, averaged these gradients, and multiplied elementwise by $\Delta$. This is a simple numerical approximation of the integral. Because it averages endpoint samples without trapezoidal endpoint weights, it is not the trapezoidal rule. Numerical completeness remains to be measured before quantitative interpretation. The initial IG example used an image of a dog.

\section{Decision-Weighted Variant}
For the decision-weighted variant, we partitioned the path into 50 coarse intervals. For interval $j$, we computed the target-logit change $F_c(\gamma((j+1)/50))-F_c(\gamma(j/50))$ and rounded that interval's fraction of the total endpoint logit change into a requested number of samples. We skipped intervals assigned zero or negative samples. At each retained point, we computed a gradient $g=\nabla_xF_c(\gamma(\alpha))$ and a directional-derivative weight $w=g\cdot\Delta$. This gives greater nominal emphasis to parts of the path where the class logit changes rapidly.

We accumulated terms proportional to $g_iw$, scaled by each interval's sample count, multiplied by $\Delta_i$, and normalized by the number of retained steps. This design is inspired by IDG rather than a numerically identical reproduction: the additional normalization and rounded, potentially skipped intervals can change the estimator relative to the published formulation \cite{idg}. A near-zero endpoint logit change can also destabilize the allocation ratio. We did not perform step-count convergence or completeness tests.

\balance
\section{Qualitative Result and Discussion}
In a macaque example, the IG visualization contains fine texture and peripheral activity, while the decision-weighted visualization appears more focused around the animal's face and body. This visual difference does not measure explanation quality. The initial IG and decision-weighted runs used different default inputs, so the example is not a controlled method comparison on the same image and target class.

A stronger experiment would process the same ImageNet validation images, classes, normalization, baseline, and color scale with both methods. Insertion and deletion curves, sensitivity to baseline and step count, runtime, and completeness error would then provide quantitative context. Reporting uncertainty across images would help distinguish a repeatable effect from a compelling single example. Heatmap concentration alone could even hide relevant but spatially distributed evidence, so it should not be the sole criterion.

\section{Conclusion}
These experiments implement path-based attribution for ResNet-50 and illustrate how decision-sensitive weighting can change a heatmap. The observed difference remains exploratory because the example is qualitative, the initial inputs differ, and the decision-weighted estimator needs numerical validation. A controlled same-image evaluation with explicit faithfulness metrics is needed to establish whether the weighting improves attribution quality.

\begin{thebibliography}{2}
\bibitem{ig} M. Sundararajan, A. Taly, and Q. Yan, ``Axiomatic attribution for deep networks,'' in \emph{Proc. 34th Int. Conf. Machine Learning}, 2017, pp. 3319--3328.
\bibitem{idg} C. Walker, S. K. Jha, K. Chen, and R. Ewetz, ``Integrated decision gradients: Compute your attributions where the model makes its decision,'' in \emph{Proc. AAAI Conf. Artificial Intelligence}, 2024.
\end{thebibliography}
\end{document}
